\documentclass[oneside,12pt]{exam}          % ----- exam

\input{../../Latex/inputs/loading_packages.tex}			% ----- Packages utiles
\input{../../Latex/inputs/preambule_controle.tex}		% ----- Préambule feuilles de TD
\input{../../Latex/inputs/mymacros.tex}                 % ----- Macros mathématiques
\input{../../Latex/inputs/mise_en_page.tex} 		   	% ----- Macros de mise en page

\geometry{left=30pt, right=30pt, top=75pt, bottom=55pt} % ----- Marges
\setlength{\headheight}{20pt}
\setlength{\headsep}{5pt}


\begin{document}

\firstpageheader
{Binôme évalué: ........................................\\[2mm]
Binôme jury : ............................................. }
{}
{Classe : ............... \\[2mm] Date : ...............}

\footer
{}
{\footnotesize \etablissement}
{}


\vspace{-1cm}\hrulefill\raisebox{-1.9pt}{\quad\decofourleft~\decosix~\decofourright \quad}\hrulefill

\begin{center}
	\textbf{\large Présentation orale évaluée}
\end{center}

%%%%%%%%%%%%%%%%%%%%%%%%%%%%%%%%%%%%%%%%%%%%%%%%%%%%%%%%%%%%%%%%%%%%%%%%%%%%%%%%%%%%%%%%%%%%%%%%%%%%%%%%%%%%%
%%%%%%%%%%%%%%%%%%%%%%%%%%%%%%%%%%%%%%%%%%%%%%%%%%%%%%%%%%%%%%%%%%%%%%%%%%%%%%%%%%%%%%%%%%%%%%%%%%%%%%%%%%%%%
%%%%%%%%%%%%%%%%%%%%%%%%%%%%%%%%%%%%%%%%%%%%%%%%%%%%%%%%%%%%%%%%%%%%%%%%%%%%%%%%%%%%%%%%%%%%%%%%%%%%%%%%%%%%%

\textbf{Entourez de manière visible les cases correspondant à vos impressions pendant l'exposé :}
{\small
\begin{center}
	\renewcommand{\arraystretch}{2}
	\begin{tabular}{| P{3.5cm} | P{3.5cm} | P{3.5cm} | P{3.5cm} | P{3.5cm} |} \hline
		\rowcolor{black} \color{white} \textbf{Qualité orale} & \color{white} \textbf{Prise de parole en continu} & \color{white} \textbf{Qualité des connaissances} & \color{white} \textbf{Qualité de l'interaction} & \color{white} \textbf{Construction de l'argumentation} \\ \hline
		\rowcolor{black!30}
		Volume sonore faible, le binôme ne capte pas l'attention. &
		Phrases courtes, longues pauses, répartition inégale du temps de parole. &
		Connaissances imprécises ou incorrectes. Difficulté à répondre aux questions. &
		Réponses courtes et ne dépassant jamais le cadre de la question. &
		Pas de compréhension du sujet, discours non argumenté et décousu. \\ \hline
		\rowcolor{black!20}
		Discours audible mais monocorde. Vocabulaire limité ou imprécis. &
		Transitions abruptes, le temps d'exposé n'est pas pleinement exploité. &
		Connaissances correctes mais difficultés à les restituer. &
		Réponses articulées, les deux élèves répondent. Pas ou peu de dépassement. &
		Le discours manque de structure, mais on peut suivre un début de raisonnement. \\ \hline
		\rowcolor{black!10}
		Quelques variations de voix, lexique adapté, le binôme capte l'attention. &
		Discours bien articulé, temps de parole égal, transitions fluides. &
		Connaissances précises, investies en répondant au jury. &
		Le binôme répond rapidement, se reprend, s'aide. &
		Des arguments précis et pertinents, qui soutiennent un raisonnement clair. \\ \hline
		La voix porte le discours, le binôme s'engage dans sa parole. Vocabulaire riche. &
		Discours fluide et efficace, les prises de paroles de l'un soutiennent celles de l'autre. &
		Connaissances maîtrisées et clairement exposées. &
		Prises d'initiative, le binôme se soutient et sait exploiter les questions posées. &
		Les enjeux sont maîtrisés, les arguments s'enchaînent de manière cohérente. \\ \hline
	\end{tabular}
\end{center}
}

\bigskip

\textbf{Écrivez ici vos questions pour le binôme qui présente l'exposé, même si vous ne les avez pas posées pendant la séance :}

\begin{enumerate}[itemsep = 5mm]
	\item .......................................................................................................................................................
	\item .......................................................................................................................................................
\end{enumerate}

\bigskip

\begin{center}
	\renewcommand{\arraystretch}{1.5}
	\begin{tabular}{| P{0.45\textwidth} | P{0.45\textwidth} | } \hline
		\textbf{Éléments positifs} & \textbf{Axes d'amélioration} \\ \hline
		\begin{itemize}[label = $\bullet$ , itemsep = 8mm]
			\item
			\item
			\item
		\end{itemize} \vspace{5mm}
	&
		\begin{itemize}[label = $\bullet$ , itemsep = 8mm]
			\item
			\item
			\item
		\end{itemize} \vspace{5mm} \\ \hline
	\end{tabular} 
\end{center}

	


\end{document}